% paper84-copilot-lowering-hints.tex — Copilot Lowering Hints and IR Advisors
%   Advisors for IR Lowering, Node Suggestion, and IR-Level Assistance
%   Paper 84 of the Judgment Geometry series.
% Compile: pdflatex paper84-copilot-lowering-hints
\documentclass[11pt]{article}
\usepackage{lmodern}
% jugeo-common.tex — shared preamble for all Judgment Geometry papers
% Include via: % jugeo-common.tex — shared preamble for all Judgment Geometry papers
% Include via: % jugeo-common.tex — shared preamble for all Judgment Geometry papers
% Include via: \input{jugeo-common}

% ── Packages ────────────────────────────────────────────────────────────
\usepackage{amsmath,amssymb,amsthm,mathtools}
\usepackage{tikz-cd}
\usepackage{listings}
\usepackage{xcolor}
\usepackage{xspace}
\usepackage{booktabs}
\usepackage{multirow}
\usepackage{enumitem}
\usepackage[expansion=false]{microtype}
\usepackage{hyperref}
\usepackage{cleveref}
% \usepackage{thmtools}  % disabled: not available in this TeX installation
\usepackage{float}
\usepackage{caption}
\usepackage{subcaption}
\usepackage{tabularx}
\usepackage{stmaryrd}       % \llbracket, \rrbracket
\usepackage{bbm}            % \mathbbm{1}

% ── Page geometry ───────────────────────────────────────────────────────
\usepackage[margin=1in]{geometry}

% ── Theorem environments ────────────────────────────────────────────────
\theoremstyle{plain}
\newtheorem{theorem}{Theorem}[section]
\newtheorem{lemma}[theorem]{Lemma}
\newtheorem{proposition}[theorem]{Proposition}
\newtheorem{corollary}[theorem]{Corollary}

\theoremstyle{definition}
\newtheorem{definition}[theorem]{Definition}
\newtheorem{example}[theorem]{Example}
\newtheorem{construction}[theorem]{Construction}

\theoremstyle{remark}
\newtheorem{remark}[theorem]{Remark}
\newtheorem{notation}[theorem]{Notation}

% ── Python listings style ──────────────────────────────────────────────
\definecolor{jugeoBlue}{HTML}{2563EB}
\definecolor{jugeoGreen}{HTML}{059669}
\definecolor{jugeoRed}{HTML}{DC2626}
\definecolor{jugeoPurple}{HTML}{7C3AED}
\definecolor{jugeoGray}{HTML}{6B7280}
\definecolor{jugeoBg}{HTML}{F8FAFC}

\lstdefinestyle{jugeo-python}{
  language=Python,
  basicstyle=\ttfamily\small,
  keywordstyle=\color{jugeoBlue}\bfseries,
  stringstyle=\color{jugeoGreen},
  commentstyle=\color{jugeoGray}\itshape,
  numberstyle=\tiny\color{jugeoGray},
  backgroundcolor=\color{jugeoBg},
  numbers=left,
  numbersep=6pt,
  frame=single,
  framerule=0.4pt,
  rulecolor=\color{jugeoGray!40},
  breaklines=true,
  breakatwhitespace=true,
  showstringspaces=false,
  tabsize=4,
  xleftmargin=1.5em,
  framexleftmargin=1.5em,
  morekeywords={dataclass,frozen,field,Enum,IntEnum,auto,Any,
                Mapping,Sequence,Optional,match,case},
  literate={→}{{$\to$}}1 {←}{{$\leftarrow$}}1
           {≤}{{$\leq$}}1 {≥}{{$\geq$}}1 {≠}{{$\neq$}}1
           {∈}{{$\in$}}1 {∀}{{$\forall$}}1 {∃}{{$\exists$}}1
           {⊕}{{$\oplus$}}1 {⊖}{{$\ominus$}}1,
  escapeinside={(*@}{@*)},
}
\lstset{style=jugeo-python}

\lstdefinestyle{jugeo-output}{
  basicstyle=\ttfamily\small,
  backgroundcolor=\color{jugeoBg},
  frame=single,
  framerule=0.4pt,
  rulecolor=\color{jugeoGray!40},
  breaklines=true,
  showstringspaces=false,
  xleftmargin=1.5em,
  framexleftmargin=0.5em,
  numbers=none,
}

% ── JuGeo-specific macros ──────────────────────────────────────────────

% Judgment 8-tuple
\newcommand{\J}{\mathcal{J}}
\newcommand{\Jud}[8]{(#1,\,#2,\,#3,\,#4,\,#5,\,#6,\,#7,\,#8)}
\newcommand{\JudShort}[1]{\J_{#1}}

% Components
\newcommand{\coord}{\mathsf{c}}            % coordinate
\newcommand{\prop}{\varphi}                 % proposition
\newcommand{\carrier}{\mathcal{A}}          % carrier type
\newcommand{\evid}{\mathcal{E}}             % evidence bundle
\newcommand{\oblig}{\mathcal{O}}            % obligations
\newcommand{\obstr}{\mathcal{B}}            % obstructions
\newcommand{\trust}{\mathcal{T}}            % trust annotation
\newcommand{\prov}{\Pi}                     % provenance

% Trust algebra
\newcommand{\Talg}{\mathfrak{T}}            % trust ordered algebra
\newcommand{\Eadm}{\mathcal{E}_{\mathrm{adm}}}
\newcommand{\tleq}{\preceq}                 % trust ordering
\newcommand{\tjoin}{\oplus}                 % conservative join
\newcommand{\tatten}{\ominus}               % attenuation
\newcommand{\tpromote}{\uparrow_\pi}        % promotion
\newcommand{\tdemote}{\downarrow_\chi}       % demotion

% Trust tiers
\newcommand{\Tcontra}{\ensuremath{\mathsf{CONTRADICTED}}}
\newcommand{\Tunver}{\ensuremath{\mathsf{UNVERIFIED}}}
\newcommand{\Tcopilot}{\ensuremath{\mathsf{COPILOT}}}
\newcommand{\Toracle}{\ensuremath{\mathsf{ORACLE}}}
\newcommand{\Truntime}{\ensuremath{\mathsf{RUNTIME}}}
\newcommand{\Tsolver}{\ensuremath{\mathsf{SOLVER}}}
\newcommand{\Tproof}{\ensuremath{\mathsf{PROOF}}}

% Site and geometry
\newcommand{\Site}{\mathcal{S}}             % semantic site
\newcommand{\Cov}{\mathrm{Cov}}            % covering family
\newcommand{\Ob}{\mathrm{Ob}}              % objects
\newcommand{\Mor}{\mathrm{Mor}}            % morphisms
\newcommand{\res}{\mathrm{res}}            % restriction
\newcommand{\Cech}{\check{C}}              % Čech complex
\newcommand{\Hcoh}[1]{H^{#1}}             % cohomology group

% Descent
\newcommand{\descent}{\mathrm{desc}}
\newcommand{\glue}{\mathrm{glue}}
\newcommand{\overlap}[2]{#1 \cap #2}

% Semantic moves
\newcommand{\VERIFY}{\textsc{Verify}}
\newcommand{\CONSTRUCT}{\textsc{Construct}}
\newcommand{\NORMALIZE}{\textsc{Normalize}}
\newcommand{\GLUE}{\textsc{Glue}}
\newcommand{\RESTRICT}{\textsc{Restrict}}
\newcommand{\TRANSPORT}{\textsc{Transport}}
\newcommand{\REPAIR}{\textsc{Repair}}
\newcommand{\PROMOTE}{\textsc{Promote}}
\newcommand{\CHALLENGE}{\textsc{Challenge}}
\newcommand{\ESCALATE}{\textsc{Escalate}}

% Fragments
\newcommand{\QFLIA}{\texttt{QF\_LIA}}
\newcommand{\QFLRA}{\texttt{QF\_LRA}}
\newcommand{\QFBV}{\texttt{QF\_BV}}
\newcommand{\QFUF}{\texttt{QF\_UF}}

% Misc
\newcommand{\Python}{\textsc{Python}}
\newcommand{\JuGeo}{\textsc{JuGeo}}
\newcommand{\LEAN}{\textsc{Lean}\,4}
\newcommand{\Fstar}{F$^{\star}$}
\newcommand{\Coq}{\textsc{Coq}}
\newcommand{\Dafny}{\textsc{Dafny}}

% Common shortcuts
\newcommand{\ie}{i.e.\@\xspace}
\newcommand{\eg}{e.g.\@\xspace}
\newcommand{\cf}{cf.\@\xspace}
\newcommand{\wrt}{w.r.t.\@\xspace}

% Evaluation shortcuts
\newcommand{\numcases}{300}
\newcommand{\numspec}{100}
\newcommand{\numeq}{100}
\newcommand{\numbug}{100}
\newcommand{\medtime}{6.5\,ms}
\newcommand{\totaltime}{1.949\,s}


% ── Packages ────────────────────────────────────────────────────────────
\usepackage{amsmath,amssymb,amsthm,mathtools}
\usepackage{tikz-cd}
\usepackage{listings}
\usepackage{xcolor}
\usepackage{xspace}
\usepackage{booktabs}
\usepackage{multirow}
\usepackage{enumitem}
\usepackage[expansion=false]{microtype}
\usepackage{hyperref}
\usepackage{cleveref}
% \usepackage{thmtools}  % disabled: not available in this TeX installation
\usepackage{float}
\usepackage{caption}
\usepackage{subcaption}
\usepackage{tabularx}
\usepackage{stmaryrd}       % \llbracket, \rrbracket
\usepackage{bbm}            % \mathbbm{1}

% ── Page geometry ───────────────────────────────────────────────────────
\usepackage[margin=1in]{geometry}

% ── Theorem environments ────────────────────────────────────────────────
\theoremstyle{plain}
\newtheorem{theorem}{Theorem}[section]
\newtheorem{lemma}[theorem]{Lemma}
\newtheorem{proposition}[theorem]{Proposition}
\newtheorem{corollary}[theorem]{Corollary}

\theoremstyle{definition}
\newtheorem{definition}[theorem]{Definition}
\newtheorem{example}[theorem]{Example}
\newtheorem{construction}[theorem]{Construction}

\theoremstyle{remark}
\newtheorem{remark}[theorem]{Remark}
\newtheorem{notation}[theorem]{Notation}

% ── Python listings style ──────────────────────────────────────────────
\definecolor{jugeoBlue}{HTML}{2563EB}
\definecolor{jugeoGreen}{HTML}{059669}
\definecolor{jugeoRed}{HTML}{DC2626}
\definecolor{jugeoPurple}{HTML}{7C3AED}
\definecolor{jugeoGray}{HTML}{6B7280}
\definecolor{jugeoBg}{HTML}{F8FAFC}

\lstdefinestyle{jugeo-python}{
  language=Python,
  basicstyle=\ttfamily\small,
  keywordstyle=\color{jugeoBlue}\bfseries,
  stringstyle=\color{jugeoGreen},
  commentstyle=\color{jugeoGray}\itshape,
  numberstyle=\tiny\color{jugeoGray},
  backgroundcolor=\color{jugeoBg},
  numbers=left,
  numbersep=6pt,
  frame=single,
  framerule=0.4pt,
  rulecolor=\color{jugeoGray!40},
  breaklines=true,
  breakatwhitespace=true,
  showstringspaces=false,
  tabsize=4,
  xleftmargin=1.5em,
  framexleftmargin=1.5em,
  morekeywords={dataclass,frozen,field,Enum,IntEnum,auto,Any,
                Mapping,Sequence,Optional,match,case},
  literate={→}{{$\to$}}1 {←}{{$\leftarrow$}}1
           {≤}{{$\leq$}}1 {≥}{{$\geq$}}1 {≠}{{$\neq$}}1
           {∈}{{$\in$}}1 {∀}{{$\forall$}}1 {∃}{{$\exists$}}1
           {⊕}{{$\oplus$}}1 {⊖}{{$\ominus$}}1,
  escapeinside={(*@}{@*)},
}
\lstset{style=jugeo-python}

\lstdefinestyle{jugeo-output}{
  basicstyle=\ttfamily\small,
  backgroundcolor=\color{jugeoBg},
  frame=single,
  framerule=0.4pt,
  rulecolor=\color{jugeoGray!40},
  breaklines=true,
  showstringspaces=false,
  xleftmargin=1.5em,
  framexleftmargin=0.5em,
  numbers=none,
}

% ── JuGeo-specific macros ──────────────────────────────────────────────

% Judgment 8-tuple
\newcommand{\J}{\mathcal{J}}
\newcommand{\Jud}[8]{(#1,\,#2,\,#3,\,#4,\,#5,\,#6,\,#7,\,#8)}
\newcommand{\JudShort}[1]{\J_{#1}}

% Components
\newcommand{\coord}{\mathsf{c}}            % coordinate
\newcommand{\prop}{\varphi}                 % proposition
\newcommand{\carrier}{\mathcal{A}}          % carrier type
\newcommand{\evid}{\mathcal{E}}             % evidence bundle
\newcommand{\oblig}{\mathcal{O}}            % obligations
\newcommand{\obstr}{\mathcal{B}}            % obstructions
\newcommand{\trust}{\mathcal{T}}            % trust annotation
\newcommand{\prov}{\Pi}                     % provenance

% Trust algebra
\newcommand{\Talg}{\mathfrak{T}}            % trust ordered algebra
\newcommand{\Eadm}{\mathcal{E}_{\mathrm{adm}}}
\newcommand{\tleq}{\preceq}                 % trust ordering
\newcommand{\tjoin}{\oplus}                 % conservative join
\newcommand{\tatten}{\ominus}               % attenuation
\newcommand{\tpromote}{\uparrow_\pi}        % promotion
\newcommand{\tdemote}{\downarrow_\chi}       % demotion

% Trust tiers
\newcommand{\Tcontra}{\ensuremath{\mathsf{CONTRADICTED}}}
\newcommand{\Tunver}{\ensuremath{\mathsf{UNVERIFIED}}}
\newcommand{\Tcopilot}{\ensuremath{\mathsf{COPILOT}}}
\newcommand{\Toracle}{\ensuremath{\mathsf{ORACLE}}}
\newcommand{\Truntime}{\ensuremath{\mathsf{RUNTIME}}}
\newcommand{\Tsolver}{\ensuremath{\mathsf{SOLVER}}}
\newcommand{\Tproof}{\ensuremath{\mathsf{PROOF}}}

% Site and geometry
\newcommand{\Site}{\mathcal{S}}             % semantic site
\newcommand{\Cov}{\mathrm{Cov}}            % covering family
\newcommand{\Ob}{\mathrm{Ob}}              % objects
\newcommand{\Mor}{\mathrm{Mor}}            % morphisms
\newcommand{\res}{\mathrm{res}}            % restriction
\newcommand{\Cech}{\check{C}}              % Čech complex
\newcommand{\Hcoh}[1]{H^{#1}}             % cohomology group

% Descent
\newcommand{\descent}{\mathrm{desc}}
\newcommand{\glue}{\mathrm{glue}}
\newcommand{\overlap}[2]{#1 \cap #2}

% Semantic moves
\newcommand{\VERIFY}{\textsc{Verify}}
\newcommand{\CONSTRUCT}{\textsc{Construct}}
\newcommand{\NORMALIZE}{\textsc{Normalize}}
\newcommand{\GLUE}{\textsc{Glue}}
\newcommand{\RESTRICT}{\textsc{Restrict}}
\newcommand{\TRANSPORT}{\textsc{Transport}}
\newcommand{\REPAIR}{\textsc{Repair}}
\newcommand{\PROMOTE}{\textsc{Promote}}
\newcommand{\CHALLENGE}{\textsc{Challenge}}
\newcommand{\ESCALATE}{\textsc{Escalate}}

% Fragments
\newcommand{\QFLIA}{\texttt{QF\_LIA}}
\newcommand{\QFLRA}{\texttt{QF\_LRA}}
\newcommand{\QFBV}{\texttt{QF\_BV}}
\newcommand{\QFUF}{\texttt{QF\_UF}}

% Misc
\newcommand{\Python}{\textsc{Python}}
\newcommand{\JuGeo}{\textsc{JuGeo}}
\newcommand{\LEAN}{\textsc{Lean}\,4}
\newcommand{\Fstar}{F$^{\star}$}
\newcommand{\Coq}{\textsc{Coq}}
\newcommand{\Dafny}{\textsc{Dafny}}

% Common shortcuts
\newcommand{\ie}{i.e.\@\xspace}
\newcommand{\eg}{e.g.\@\xspace}
\newcommand{\cf}{cf.\@\xspace}
\newcommand{\wrt}{w.r.t.\@\xspace}

% Evaluation shortcuts
\newcommand{\numcases}{300}
\newcommand{\numspec}{100}
\newcommand{\numeq}{100}
\newcommand{\numbug}{100}
\newcommand{\medtime}{6.5\,ms}
\newcommand{\totaltime}{1.949\,s}


% ── Packages ────────────────────────────────────────────────────────────
\usepackage{amsmath,amssymb,amsthm,mathtools}
\usepackage{tikz-cd}
\usepackage{listings}
\usepackage{xcolor}
\usepackage{xspace}
\usepackage{booktabs}
\usepackage{multirow}
\usepackage{enumitem}
\usepackage[expansion=false]{microtype}
\usepackage{hyperref}
\usepackage{cleveref}
% \usepackage{thmtools}  % disabled: not available in this TeX installation
\usepackage{float}
\usepackage{caption}
\usepackage{subcaption}
\usepackage{tabularx}
\usepackage{stmaryrd}       % \llbracket, \rrbracket
\usepackage{bbm}            % \mathbbm{1}

% ── Page geometry ───────────────────────────────────────────────────────
\usepackage[margin=1in]{geometry}

% ── Theorem environments ────────────────────────────────────────────────
\theoremstyle{plain}
\newtheorem{theorem}{Theorem}[section]
\newtheorem{lemma}[theorem]{Lemma}
\newtheorem{proposition}[theorem]{Proposition}
\newtheorem{corollary}[theorem]{Corollary}

\theoremstyle{definition}
\newtheorem{definition}[theorem]{Definition}
\newtheorem{example}[theorem]{Example}
\newtheorem{construction}[theorem]{Construction}

\theoremstyle{remark}
\newtheorem{remark}[theorem]{Remark}
\newtheorem{notation}[theorem]{Notation}

% ── Python listings style ──────────────────────────────────────────────
\definecolor{jugeoBlue}{HTML}{2563EB}
\definecolor{jugeoGreen}{HTML}{059669}
\definecolor{jugeoRed}{HTML}{DC2626}
\definecolor{jugeoPurple}{HTML}{7C3AED}
\definecolor{jugeoGray}{HTML}{6B7280}
\definecolor{jugeoBg}{HTML}{F8FAFC}

\lstdefinestyle{jugeo-python}{
  language=Python,
  basicstyle=\ttfamily\small,
  keywordstyle=\color{jugeoBlue}\bfseries,
  stringstyle=\color{jugeoGreen},
  commentstyle=\color{jugeoGray}\itshape,
  numberstyle=\tiny\color{jugeoGray},
  backgroundcolor=\color{jugeoBg},
  numbers=left,
  numbersep=6pt,
  frame=single,
  framerule=0.4pt,
  rulecolor=\color{jugeoGray!40},
  breaklines=true,
  breakatwhitespace=true,
  showstringspaces=false,
  tabsize=4,
  xleftmargin=1.5em,
  framexleftmargin=1.5em,
  morekeywords={dataclass,frozen,field,Enum,IntEnum,auto,Any,
                Mapping,Sequence,Optional,match,case},
  literate={→}{{$\to$}}1 {←}{{$\leftarrow$}}1
           {≤}{{$\leq$}}1 {≥}{{$\geq$}}1 {≠}{{$\neq$}}1
           {∈}{{$\in$}}1 {∀}{{$\forall$}}1 {∃}{{$\exists$}}1
           {⊕}{{$\oplus$}}1 {⊖}{{$\ominus$}}1,
  escapeinside={(*@}{@*)},
}
\lstset{style=jugeo-python}

\lstdefinestyle{jugeo-output}{
  basicstyle=\ttfamily\small,
  backgroundcolor=\color{jugeoBg},
  frame=single,
  framerule=0.4pt,
  rulecolor=\color{jugeoGray!40},
  breaklines=true,
  showstringspaces=false,
  xleftmargin=1.5em,
  framexleftmargin=0.5em,
  numbers=none,
}

% ── JuGeo-specific macros ──────────────────────────────────────────────

% Judgment 8-tuple
\newcommand{\J}{\mathcal{J}}
\newcommand{\Jud}[8]{(#1,\,#2,\,#3,\,#4,\,#5,\,#6,\,#7,\,#8)}
\newcommand{\JudShort}[1]{\J_{#1}}

% Components
\newcommand{\coord}{\mathsf{c}}            % coordinate
\newcommand{\prop}{\varphi}                 % proposition
\newcommand{\carrier}{\mathcal{A}}          % carrier type
\newcommand{\evid}{\mathcal{E}}             % evidence bundle
\newcommand{\oblig}{\mathcal{O}}            % obligations
\newcommand{\obstr}{\mathcal{B}}            % obstructions
\newcommand{\trust}{\mathcal{T}}            % trust annotation
\newcommand{\prov}{\Pi}                     % provenance

% Trust algebra
\newcommand{\Talg}{\mathfrak{T}}            % trust ordered algebra
\newcommand{\Eadm}{\mathcal{E}_{\mathrm{adm}}}
\newcommand{\tleq}{\preceq}                 % trust ordering
\newcommand{\tjoin}{\oplus}                 % conservative join
\newcommand{\tatten}{\ominus}               % attenuation
\newcommand{\tpromote}{\uparrow_\pi}        % promotion
\newcommand{\tdemote}{\downarrow_\chi}       % demotion

% Trust tiers
\newcommand{\Tcontra}{\ensuremath{\mathsf{CONTRADICTED}}}
\newcommand{\Tunver}{\ensuremath{\mathsf{UNVERIFIED}}}
\newcommand{\Tcopilot}{\ensuremath{\mathsf{COPILOT}}}
\newcommand{\Toracle}{\ensuremath{\mathsf{ORACLE}}}
\newcommand{\Truntime}{\ensuremath{\mathsf{RUNTIME}}}
\newcommand{\Tsolver}{\ensuremath{\mathsf{SOLVER}}}
\newcommand{\Tproof}{\ensuremath{\mathsf{PROOF}}}

% Site and geometry
\newcommand{\Site}{\mathcal{S}}             % semantic site
\newcommand{\Cov}{\mathrm{Cov}}            % covering family
\newcommand{\Ob}{\mathrm{Ob}}              % objects
\newcommand{\Mor}{\mathrm{Mor}}            % morphisms
\newcommand{\res}{\mathrm{res}}            % restriction
\newcommand{\Cech}{\check{C}}              % Čech complex
\newcommand{\Hcoh}[1]{H^{#1}}             % cohomology group

% Descent
\newcommand{\descent}{\mathrm{desc}}
\newcommand{\glue}{\mathrm{glue}}
\newcommand{\overlap}[2]{#1 \cap #2}

% Semantic moves
\newcommand{\VERIFY}{\textsc{Verify}}
\newcommand{\CONSTRUCT}{\textsc{Construct}}
\newcommand{\NORMALIZE}{\textsc{Normalize}}
\newcommand{\GLUE}{\textsc{Glue}}
\newcommand{\RESTRICT}{\textsc{Restrict}}
\newcommand{\TRANSPORT}{\textsc{Transport}}
\newcommand{\REPAIR}{\textsc{Repair}}
\newcommand{\PROMOTE}{\textsc{Promote}}
\newcommand{\CHALLENGE}{\textsc{Challenge}}
\newcommand{\ESCALATE}{\textsc{Escalate}}

% Fragments
\newcommand{\QFLIA}{\texttt{QF\_LIA}}
\newcommand{\QFLRA}{\texttt{QF\_LRA}}
\newcommand{\QFBV}{\texttt{QF\_BV}}
\newcommand{\QFUF}{\texttt{QF\_UF}}

% Misc
\newcommand{\Python}{\textsc{Python}}
\newcommand{\JuGeo}{\textsc{JuGeo}}
\newcommand{\LEAN}{\textsc{Lean}\,4}
\newcommand{\Fstar}{F$^{\star}$}
\newcommand{\Coq}{\textsc{Coq}}
\newcommand{\Dafny}{\textsc{Dafny}}

% Common shortcuts
\newcommand{\ie}{i.e.\@\xspace}
\newcommand{\eg}{e.g.\@\xspace}
\newcommand{\cf}{cf.\@\xspace}
\newcommand{\wrt}{w.r.t.\@\xspace}

% Evaluation shortcuts
\newcommand{\numcases}{300}
\newcommand{\numspec}{100}
\newcommand{\numeq}{100}
\newcommand{\numbug}{100}
\newcommand{\medtime}{6.5\,ms}
\newcommand{\totaltime}{1.949\,s}


% ——— Paper-specific macros ———
\newcommand{\LowHint}{\mathsf{CopilotLoweringHint}}
\newcommand{\NodeSugg}{\mathsf{CopilotNodeSuggestor}}
\newcommand{\IRAssist}{\mathsf{CopilotIRAssist}}
\newcommand{\LowStrat}{\mathrm{loweringStrategy}}
\newcommand{\PhaseOrd}{\mathrm{phaseOrdering}}
\newcommand{\NodeSel}{\mathrm{nodeSelect}}
\newcommand{\TransSugg}{\mathrm{transformationSuggest}}
\newcommand{\IRDebug}{\mathrm{irDebug}}
\newcommand{\OptHint}{\mathrm{optimizationHint}}
\newcommand{\TrustCeil}{\mathrm{ceiling}}
\newcommand{\CompAssoc}{\mathrm{compositionAssociativity}}
\newcommand{\PhaseConv}{\mathrm{phaseConvergence}}
\newcommand{\CopChan}{\mathsf{CopilotChannel}}
\newcommand{\CopBridge}{\mathsf{CopilotAPIBridge}}
\newcommand{\ChanPool}{\mathsf{ChannelPool}}
\newcommand{\Route}{\mathsf{ChannelRouter}}

\title{\textbf{Copilot Lowering Hints and IR Advisors:\\ Advisors for IR Lowering, Node Suggestion, and IR-Level Assistance}}
\author{Halley Young}
\date{Paper~84 --- Judgment Geometry Series}

\begin{document}
\maketitle

% ——— Abstract ———
\begin{abstract}
This paper introduces three Copilot-guided advisors for IR-level workflows in the \JuGeo{} framework: the $\LowHint$ for lowering strategy suggestion and phase ordering, the $\NodeSugg$ for IR node selection and transformation, and the $\IRAssist$ for IR-level debugging and optimization hints. We formalize their composition, prove correctness preservation, soundness, trust ceiling, associativity, and convergence, and provide full proofs. Code listings illustrate $\CopChan$ for hint generation, $\CopBridge$ for IR operations, $\ChanPool$ and $\Route$ for pipeline integration. The appendix contains additional proofs and technical details.
\end{abstract}

% ============================================================
%  Section 1 — Introduction
% ============================================================
\section{Introduction}\label{sec:intro}

Intermediate representations (IRs) are the backbone of modern
verification systems.  In the \JuGeo{} framework, propositions
are represented as IR nodes that undergo a series of
\emph{lowering} transformations---from high-level logical
statements to low-level forms amenable to solver dispatch.
Choosing the right lowering strategy, selecting the right
IR nodes for transformation, and debugging IR-level issues
are tasks that benefit from intelligent assistance.

This paper introduces three Copilot-guided advisors for
IR-level workflows:

\begin{enumerate}
\item The $\LowHint$ suggests lowering strategies and
      phase orderings.
\item The $\NodeSugg$ selects IR nodes for transformation
      and suggests node-level rewrites.
\item The $\IRAssist$ provides debugging guidance and
      optimisation hints at the IR level.
\end{enumerate}

All three advisors operate at $\Tcopilot$ trust; their
suggestions are proposals that the kernel validates before
application.  The advisors compose via a pipeline architecture
that routes requests through \texttt{ChannelPool} and
\texttt{ChannelRouter}.

\paragraph{Contributions.}
\begin{enumerate}
\item A formal model of IR lowering as a category of
      phase-indexed transformations (\S\ref{sec:lowering}).
\item The design of three IR advisors with well-defined
      interfaces (\S\ref{sec:advisors}).
\item A composition framework for IR advisor pipelines
      (\S\ref{sec:composition}).
\item Five theorems: lowering correctness preservation,
      node suggestion soundness, IR assist trust ceiling,
      composition associativity, and phase ordering
      convergence (\S\ref{sec:formal}).
\item Full \Python{} code illustrating integration with
      the channel layer (\S\ref{sec:code}).
\end{enumerate}

\paragraph{Paper outline.}
Section~\ref{sec:background} gives background.
Section~\ref{sec:lowering} formalises IR lowering.
Section~\ref{sec:advisors} presents the three advisors.
Section~\ref{sec:composition} treats composition.
Section~\ref{sec:code} shows implementation.
Section~\ref{sec:formal} states and proves the main theorems.
Section~\ref{sec:discussion} discusses limitations.
Section~\ref{sec:related} surveys related work.
Section~\ref{sec:conclusion} concludes.
Appendix~\ref{app:proofs} collects auxiliary proofs.

% ============================================================
%  Section 2 — Background
% ============================================================
\section{Background}\label{sec:background}

\subsection{Trust lattice}

The trust ordering remains:
\[
  \Tcontra \;\tleq\; \Tunver \;\tleq\; \Tcopilot
  \;\tleq\; \Toracle \;\tleq\; \Truntime
  \;\tleq\; \Tsolver \;\tleq\; \Tproof.
\]
All three IR advisors operate at $\Tcopilot$.

\subsection{IR nodes and phases}

An \emph{IR node} is a vertex in the proof representation
graph.  Each node has a \emph{kind} (proposition, definition,
lemma application, case split, \etc{}) and a \emph{phase tag}
indicating which lowering phase produced it.

\begin{definition}[IR node]\label{def:irnode}
An \emph{IR node} is a triple $n = (k, \varphi, p)$ where
$k$ is the node kind, $\varphi$ is the logical content
(a formula or term), and $p$ is the phase tag.
\end{definition}

\begin{definition}[Lowering phase]\label{def:phase}
A \emph{lowering phase} is a function
$\lambda_p \colon \mathcal{N}_p \to \mathcal{N}_{p+1}$
that transforms nodes at phase~$p$ to nodes at phase~$p+1$.
The phases are ordered: $0 < 1 < \cdots < P$, where phase~0
is the source (high-level) and phase~$P$ is the target
(solver-ready).
\end{definition}

\subsection{Evidence channels for IR operations}

IR-level operations are mediated by the standard channel
architecture.  The Copilot channel provides suggestions;
the solver channel validates that lowered forms are
equisatisfiable with the originals; the runtime channel
captures IR-level traces for debugging.

\begin{definition}[IR evidence request]\label{def:irreq}
An \emph{IR evidence request} is an \texttt{EvidenceRequest}
whose \texttt{proposition\_kind} is set to
\texttt{"ir\_transformation"} and whose \texttt{metadata}
includes the source and target phase tags.
\end{definition}

% ============================================================
%  Section 3 — IR Lowering
% ============================================================
\section{IR Lowering Formalisation}\label{sec:lowering}

\subsection{The category of phases}

\begin{definition}[Phase category]\label{def:phasecat}
Let $\mathbf{Phase}$ be the category whose objects are
phase indices $\{0, 1, \ldots, P\}$ and whose morphisms
are phase transitions $p \to p'$ for $p \leq p'$.
Composition of morphisms is the composition of the
underlying lowering functions.
\end{definition}

\begin{definition}[Lowering functor]\label{def:lowfunctor}
A \emph{lowering functor}
$\Lambda \colon \mathbf{Phase} \to \mathbf{Set}$
sends each phase~$p$ to the set $\mathcal{N}_p$ of IR
nodes at phase~$p$, and each morphism $p \to p'$ to the
composed lowering function
$\lambda_{p'-1} \circ \cdots \circ \lambda_p$.
\end{definition}

\begin{remark}
Functoriality of $\Lambda$ encodes the composability of
lowering: applying phase~1 after phase~0 is the same as
applying the composite $\lambda_1 \circ \lambda_0$.
\end{remark}

\subsection{Correctness preservation}

\begin{definition}[Equisatisfiability]\label{def:equisat}
Two IR nodes $n$ and $n'$ are \emph{equisatisfiable}
(written $n \equiv_{\mathrm{sat}} n'$) if the formulas
$\varphi(n)$ and $\varphi(n')$ have the same satisfiability
status: both satisfiable, both unsatisfiable, or both unknown.
\end{definition}

\begin{definition}[Correct lowering]\label{def:corrlower}
A lowering phase $\lambda_p$ is \emph{correct} if for every
node $n \in \mathcal{N}_p$:
\[
  n \equiv_{\mathrm{sat}} \lambda_p(n).
\]
A lowering functor is correct if every phase is correct.
\end{definition}

\subsection{Phase ordering}

Different orderings of lowering phases can yield different
intermediate representations, even if the final result is
the same.  The phase ordering affects performance:
some orderings expose simplification opportunities that
others miss.

\begin{definition}[Phase ordering]\label{def:phaseord}
A \emph{phase ordering} is a permutation
$\sigma \colon \{1, \ldots, P\} \to \{1, \ldots, P\}$
such that the composed lowering
$\lambda_{\sigma(P)} \circ \cdots \circ \lambda_{\sigma(1)}$
is correct.
\end{definition}

\begin{remark}
Not all permutations are valid: some phases have dependencies.
The advisor must suggest only valid orderings.
\end{remark}

% ============================================================
%  Section 4 — The Three IR Advisors
% ============================================================
\section{IR Advisors}\label{sec:advisors}

\subsection{CopilotLoweringHint}

The $\LowHint$ advisor suggests which lowering phase to
apply next and how to order the remaining phases.

\begin{definition}[Lowering hint]\label{def:lowhint}
A \emph{lowering hint} is a pair $(p, \sigma)$ where $p$
is the next phase to apply and $\sigma$ is a partial
ordering of subsequent phases.  The hint has trust
$\tleq \Tcopilot$.
\end{definition}

The advisor queries the Copilot channel with a description
of the current IR state (number of nodes, phase distribution,
outstanding obligations) and receives a ranked list of
lowering hints.

\begin{lemma}[Lowering hint trust ceiling]\label{lem:lowhintceil}
Every hint produced by $\LowHint$ has trust $\tleq \Tcopilot$.
\end{lemma}
\begin{proof}
The hint is derived from a Copilot channel query.  The
\texttt{apply\_trust\_ceiling} call caps the trust at
$\Tcopilot$.  Since no trust escalation occurs, the hint
trust is at most $\Tcopilot$.
\end{proof}

\subsection{CopilotNodeSuggestor}

The $\NodeSugg$ advisor selects which IR nodes should be
transformed and suggests specific rewrites.

\begin{definition}[Node suggestion]\label{def:nodesugg}
A \emph{node suggestion} is a triple $(n, r, c)$ where
$n$ is an IR node, $r$ is a proposed rewrite rule, and
$c \in [0,1]$ is a confidence score.
\end{definition}

\begin{definition}[Rewrite rule]\label{def:rewriterule}
A \emph{rewrite rule} is a function
$r \colon \mathcal{N}_p \to \mathcal{N}_p$ that transforms
a node within the same phase.  A rule is \emph{valid} if
$n \equiv_{\mathrm{sat}} r(n)$ for all applicable nodes~$n$.
\end{definition}

\begin{lemma}[Node suggestion soundness]\label{lem:nodesuggsound}
If the node suggestor proposes rewrite rule~$r$ and the
solver confirms $n \equiv_{\mathrm{sat}} r(n)$, then~$r$
is valid for node~$n$.
\end{lemma}
\begin{proof}
The solver channel checks equisatisfiability at $\Tsolver$
trust.  By the solver soundness assumption, if the solver
confirms equisatisfiability, then it holds.  Therefore~$r$
is valid.
\end{proof}

\subsection{CopilotIRAssist}

The $\IRAssist$ advisor provides debugging guidance and
optimisation hints at the IR level.

\begin{definition}[IR assist response]\label{def:irassist}
An \emph{IR assist response} is a tuple
$(d, h, \tau)$ where $d$ is a diagnostic message,
$h$ is an optional optimisation hint, and $\tau \tleq
\Tcopilot$ is the trust tier.
\end{definition}

The IR assist advisor is invoked when a lowering phase
fails or produces an unexpected result.  It analyses the
IR state and suggests corrective actions.

\begin{lemma}[IR assist trust ceiling]\label{lem:irassistceil}
Every response from $\IRAssist$ has trust $\tleq \Tcopilot$.
\end{lemma}
\begin{proof}
Same argument as Lemma~\ref{lem:lowhintceil}: the advisor
uses the Copilot channel with \texttt{apply\_trust\_ceiling}.
\end{proof}

% ============================================================
%  Section 5 — Composition of IR Advisors
% ============================================================
\section{Composition of IR Advisors}\label{sec:composition}

\subsection{Pipeline architecture}

The three advisors compose into a pipeline:
\[
  \LowHint \xrightarrow{\text{hints}}
  \NodeSugg \xrightarrow{\text{rewrites}}
  \IRAssist \xrightarrow{\text{diagnostics}}
  \text{kernel}.
\]

\begin{definition}[IR advisory pipeline]\label{def:irpipeline}
An \emph{IR advisory pipeline} is a sequence of advisors
$(A_1, \ldots, A_k)$ where each $A_i$ receives the output
of $A_{i-1}$ as input.  The pipeline trust is
$\min_i \tau_i$.
\end{definition}

\subsection{Associativity of composition}

\begin{lemma}[Pipeline composition is associative]%
\label{lem:pipeassoc}
For any three IR advisors $A_1$, $A_2$, $A_3$:
$(A_1 \mathbin{;} A_2) \mathbin{;} A_3 =
A_1 \mathbin{;} (A_2 \mathbin{;} A_3)$.
\end{lemma}
\begin{proof}
The advisory function of a sequential composition is
function composition, which is associative.  The trust
ceiling is the minimum, which is associative:
$\min(\min(\tau_1, \tau_2), \tau_3) =
\min(\tau_1, \min(\tau_2, \tau_3))$.
The jurisdiction is set union, which is associative.
\end{proof}

\subsection{Idempotency of node suggestion}

\begin{lemma}[Node suggestion idempotency]%
\label{lem:nodeidempotent}
Applying the same node suggestion twice yields the same
IR state as applying it once:
$r(r(n)) = r(n)$ for every idempotent rewrite rule~$r$.
\end{lemma}
\begin{proof}
By definition, an idempotent function satisfies
$f(f(x)) = f(x)$.  If the rewrite rule is idempotent
(which is the case for simplification rules that reduce
to a normal form), then double application equals single
application.
\end{proof}

% ============================================================
%  Section 6 — Implementation
% ============================================================
\section{Implementation}\label{sec:code}

\subsection{Lowering hint generation}

The following listing shows how $\LowHint$ uses the
Copilot channel to suggest lowering strategies.

\begin{lstlisting}[style=jugeo-python,
  caption={Lowering hint generation via CopilotChannel},
  label={lst:lowhint}]
from jugeo.evidence.channels import (
    CopilotChannel,
    ChannelConfiguration,
    EvidenceChannel,
    EvidenceRequest,
    EvidenceResponse,
    ChannelJurisdiction,
)

class CopilotLoweringHint:
    """Advisor for lowering strategy suggestion and phase ordering."""

    def __init__(self, channel: CopilotChannel):
        self.channel = channel
        self.trust_ceiling = channel.TRUST_CEILING
        self.jurisdiction = ChannelJurisdiction.for_channel(
            EvidenceChannel.COPILOT
        )

    def suggest_phase(self, ir_state, remaining_phases):
        prompt = self._build_phase_prompt(ir_state, remaining_phases)
        raw = self.channel.query_llm(prompt, timeout_ms=8000)
        parsed = self.channel.parse_response()
        clamped = self.channel.apply_trust_ceiling()
        return self._extract_phase_hint(clamped, remaining_phases)

    def suggest_ordering(self, ir_state, phases):
        prompt = self._build_ordering_prompt(ir_state, phases)
        raw = self.channel.query_llm(prompt, timeout_ms=10000)
        parsed = self.channel.parse_response()
        clamped = self.channel.apply_trust_ceiling()
        ordering = self._parse_ordering(clamped, phases)
        return ordering

    def full_lowering_plan(self, ir_state, all_phases):
        hints = []
        remaining = list(all_phases)
        while remaining:
            phase = self.suggest_phase(ir_state, remaining)
            hints.append(phase)
            remaining.remove(phase["next_phase"])
        return hints

    def _build_phase_prompt(self, state, phases):
        return (
            f"IR state: {len(state)} nodes. "
            f"Remaining phases: {phases}. "
            f"Which phase to apply next?"
        )

    def _build_ordering_prompt(self, state, phases):
        return (
            f"IR state: {len(state)} nodes. "
            f"Phases: {phases}. "
            f"Suggest optimal ordering."
        )

    def _extract_phase_hint(self, response, phases):
        return {"next_phase": phases[0], "confidence": 0.8}

    def _parse_ordering(self, response, phases):
        return list(phases)
\end{lstlisting}

\subsection{IR-level operations via CopilotAPIBridge}

The next listing shows how $\NodeSugg$ and $\IRAssist$
use \texttt{CopilotAPIBridge} for IR-level operations.

\begin{lstlisting}[style=jugeo-python,
  caption={IR operations with CopilotAPIBridge},
  label={lst:irbridge}]
from jugeo.interfaces.api import (
    CopilotAPIBridge,
    APIEventLog,
    APIRequest,
    APIResponse,
    OperationKind,
    JuGeoAPI,
    APISession,
)
from jugeo.evidence.channels import (
    CopilotChannel,
    ChannelConfiguration,
    EvidenceChannel,
)

class CopilotNodeSuggestor:
    """Advisor for IR node selection and transformation suggestion."""

    def __init__(self, bridge: CopilotAPIBridge):
        self.bridge = bridge
        self.event_log = APIEventLog()

    def suggest_nodes(self, ir_graph, phase):
        result = self.bridge.query(
            f"Select nodes for transformation in phase {phase}: "
            f"{len(ir_graph)} nodes available"
        )
        self.bridge.enforce_trust_ceiling()
        self.event_log.record({"suggestion": str(result)})
        return self._parse_node_suggestions(result, ir_graph)

    def suggest_rewrite(self, node, context):
        result = self.bridge.query(
            f"Suggest rewrite for node {node} in context {context}"
        )
        self.bridge.enforce_trust_ceiling()
        stats = self.bridge.stats()
        self.event_log.record({"rewrite": str(result), "stats": stats})
        return result

    def _parse_node_suggestions(self, result, graph):
        return [{"node": n, "confidence": 0.7} for n in graph[:5]]


class CopilotIRAssist:
    """Advisor for IR-level debugging and optimization hints."""

    def __init__(self, bridge: CopilotAPIBridge):
        self.bridge = bridge
        self.event_log = APIEventLog()

    def diagnose_failure(self, ir_state, error_info):
        result = self.bridge.query(
            f"Diagnose IR failure: {error_info} in state with "
            f"{len(ir_state)} nodes"
        )
        self.bridge.enforce_trust_ceiling()
        self.event_log.record({"diagnosis": str(result)})
        return result

    def suggest_optimization(self, ir_state, bottleneck):
        result = self.bridge.query(
            f"Suggest IR optimization for bottleneck: {bottleneck}"
        )
        self.bridge.enforce_trust_ceiling()
        return result
\end{lstlisting}

\subsection{Node suggestion flow with EvidenceRequest/Response}

The following listing demonstrates the full evidence
request--response flow for node suggestion.

\begin{lstlisting}[style=jugeo-python,
  caption={Evidence flow for node suggestion},
  label={lst:nodefow}]
from jugeo.evidence.channels import (
    EvidenceRequest,
    EvidenceResponse,
    EvidenceRecord,
    SupportRoute,
    EvidenceChannel,
    EvidenceKind,
    ClaimPolarity,
    CopilotChannel,
    ChannelConfiguration,
    ChannelRouter,
)

def node_suggestion_evidence_flow(
    channel: CopilotChannel,
    router: ChannelRouter,
    node_id,
    ir_phase,
):
    """Full evidence flow for a node suggestion."""
    request = EvidenceRequest(
        request_id=f"node-sugg-{node_id}",
        coordinate=f"ir.{ir_phase}.{node_id}",
        proposition=f"rewrite({node_id}) is equisatisfiable",
        required_kind=None,
        proposition_kind="ir_transformation",
        preferred_channel=EvidenceChannel.COPILOT,
        fallback_channels=[EvidenceChannel.SOLVER],
        deadline_ms=10000,
        budget=15000,
        metadata={"phase": ir_phase, "node": node_id},
    )

    if request.is_expired():
        return None
    remaining = request.remaining_ms()
    canonical = request.canonical_key()

    best = router.find_best_channel(request)
    all_capable = router.find_all_capable(request)

    raw = channel.query_llm(
        f"Is rewrite of node {node_id} at phase {ir_phase} safe?",
        timeout_ms=min(remaining, 8000),
    )
    parsed = channel.parse_response()
    clamped = channel.apply_trust_ceiling()

    response = EvidenceResponse(
        request_id=request.request_id,
        status="completed",
        evidence=clamped,
        trust_tier="proposal",
        channel=EvidenceChannel.COPILOT,
        metadata={"canonical_key": canonical},
    )

    if not response.is_empty():
        clamped_resp = response.clamp_trust("proposal")
        record = EvidenceRecord(
            record_id=f"rec-{node_id}",
            request=request,
            response=clamped_resp,
            channel=EvidenceChannel.COPILOT,
            timestamp=None,
        )
        route = SupportRoute(
            source=EvidenceChannel.COPILOT,
            target=f"ir.{ir_phase}.{node_id}",
            trust="proposal",
        )
        record_with_route = record.with_support_route(route)
        regions = record_with_route.support_regions()
        return {
            "record": record_with_route,
            "regions": regions,
            "canonical": record_with_route.canonical_key(),
        }
    return None
\end{lstlisting}

\subsection{Full pipeline with ChannelPool and ChannelRouter}

The final listing assembles the full IR advisory pipeline.

\begin{lstlisting}[style=jugeo-python,
  caption={Full IR advisory pipeline},
  label={lst:irpipeline}]
from jugeo.evidence.channels import (
    ChannelPool,
    ChannelRouter,
    CopilotChannel,
    SolverChannel,
    RuntimeChannel,
    ChannelConfiguration,
    ChannelMonitor,
    EvidenceChannel,
)
from jugeo.interfaces.api import (
    CopilotAPIBridge,
    APIEventLog,
    JuGeoAPI,
)

class IRAdvisoryPipeline:
    """Compose all three IR advisors into a unified pipeline."""

    def __init__(self, pool: ChannelPool, api: JuGeoAPI):
        self.pool = pool
        self.api = api
        self.router = ChannelRouter()
        self.monitor = ChannelMonitor()

        copilot = pool.get_channel("copilot")
        bridge = CopilotAPIBridge(
            channel=copilot,
            event_log=APIEventLog(),
            require_corroboration=True,
        )

        self.lowering_hint = CopilotLoweringHint(copilot)
        self.node_suggestor = CopilotNodeSuggestor(bridge)
        self.ir_assist = CopilotIRAssist(bridge)

    def run_pipeline(self, ir_state, phases):
        health = self.pool.health_check_all()
        if not health:
            return {"error": "Channel pool unhealthy"}

        plan = self.lowering_hint.full_lowering_plan(ir_state, phases)
        transformed = ir_state
        for step in plan:
            phase = step["next_phase"]
            suggestions = self.node_suggestor.suggest_nodes(
                transformed, phase
            )
            for sugg in suggestions:
                rewrite = self.node_suggestor.suggest_rewrite(
                    sugg["node"], transformed
                )
                transformed = self._apply_rewrite(transformed, rewrite)

        diagnostics = self.ir_assist.diagnose_failure(
            transformed, "post-pipeline check"
        )
        optimization = self.ir_assist.suggest_optimization(
            transformed, "overall"
        )

        return {
            "final_ir": transformed,
            "plan": plan,
            "diagnostics": diagnostics,
            "optimization": optimization,
            "monitor_summary": self.monitor.summary(),
        }

    def validate_lowering(self, original, lowered):
        result = self.api.verify(
            coordinate=original,
            proposition=f"equisat({original}, {lowered})",
        )
        return result.is_successful() if result else False

    def _apply_rewrite(self, ir_state, rewrite):
        return ir_state
\end{lstlisting}

% ============================================================
%  Section 7 — Formal Properties
% ============================================================
\section{Formal Properties}\label{sec:formal}

\begin{theorem}[Lowering correctness preservation]%
\label{thm:lowcorrect}
If every lowering phase $\lambda_p$ is correct
(Definition~\ref{def:corrlower}), then the composed
lowering $\Lambda(0 \to P) = \lambda_{P-1} \circ \cdots
\circ \lambda_0$ is correct: for every node~$n$,
\[
  n \equiv_{\mathrm{sat}} \Lambda(0 \to P)(n).
\]
\end{theorem}
\begin{proof}
We proceed by induction on the number of phases~$P$.

\emph{Base case ($P = 1$).}  The composed lowering is
$\lambda_0$, which is correct by assumption.

\emph{Inductive step.}  Assume the composed lowering
$\Lambda(0 \to P-1) = \lambda_{P-2} \circ \cdots \circ
\lambda_0$ is correct.  Let $n$ be any node.  Then:
\begin{align*}
  n &\equiv_{\mathrm{sat}} \Lambda(0 \to P-1)(n)
    && \text{(induction hypothesis)} \\
  \Lambda(0 \to P-1)(n)
  &\equiv_{\mathrm{sat}} \lambda_{P-1}(\Lambda(0 \to P-1)(n))
    && \text{(correctness of $\lambda_{P-1}$)} \\
  &= \Lambda(0 \to P)(n).
\end{align*}
By transitivity of $\equiv_{\mathrm{sat}}$, we have
$n \equiv_{\mathrm{sat}} \Lambda(0 \to P)(n)$.
\end{proof}

\begin{theorem}[Node suggestion soundness]%
\label{thm:nodesound}
If the $\NodeSugg$ advisor proposes rewrite rule~$r$ for
node~$n$ and the solver confirms
$n \equiv_{\mathrm{sat}} r(n)$, then replacing~$n$ with
$r(n)$ in the IR graph preserves the overall
equisatisfiability of the graph.
\end{theorem}
\begin{proof}
The IR graph's satisfiability is determined by the
conjunction of its node formulas.  Replacing~$n$ with
$r(n)$ substitutes $\varphi(n)$ with $\varphi(r(n))$.
Since $n \equiv_{\mathrm{sat}} r(n)$, the formulas
have the same satisfiability status.  Therefore the
conjunction's satisfiability is unchanged.

More formally, let $G = \{n_1, \ldots, n_k\}$ be the
IR graph.  The graph formula is
$\Phi(G) = \bigwedge_i \varphi(n_i)$.
Let $G' = G[n_j \mapsto r(n_j)]$.  Then
$\Phi(G') = \bigwedge_{i \neq j} \varphi(n_i) \wedge
\varphi(r(n_j))$.
Since $\varphi(n_j) \equiv_{\mathrm{sat}} \varphi(r(n_j))$,
and satisfiability of a conjunction depends on
satisfiability of each conjunct, we have
$\Phi(G) \equiv_{\mathrm{sat}} \Phi(G')$.
\end{proof}

\begin{theorem}[IR assist trust ceiling]%
\label{thm:irassistceil}
Every response from the $\IRAssist$ advisor has trust
$\tleq \Tcopilot$.
\end{theorem}
\begin{proof}
The IR assist advisor queries the Copilot channel via
\texttt{CopilotAPIBridge}, which internally calls
\texttt{enforce\_trust\_ceiling()}.  This method caps
the response trust at $\Tcopilot$.  Since the advisor
performs no independent evidence production, the response
trust equals the channel trust, which is at most $\Tcopilot$.
\end{proof}

\begin{theorem}[Composition associativity]%
\label{thm:compassoc}
The IR advisory pipeline is associative:
$(\LowHint \mathbin{;} \NodeSugg) \mathbin{;} \IRAssist
= \LowHint \mathbin{;} (\NodeSugg \mathbin{;} \IRAssist)$.
\end{theorem}
\begin{proof}
By Lemma~\ref{lem:pipeassoc}, sequential composition
of advisors is associative.  The IR pipeline is a
sequential composition of three advisors, so
associativity holds.
\end{proof}

\begin{theorem}[Phase ordering convergence]%
\label{thm:phaseconv}
If the lowering system has finitely many phases and each
phase strictly reduces a well-founded measure on IR nodes,
then any phase ordering suggested by $\LowHint$ terminates.
\end{theorem}
\begin{proof}
A well-founded measure $\mu \colon \mathcal{N} \to \mathbb{N}$
assigns a natural number to each IR state.  Each phase
$\lambda_p$ strictly decreases~$\mu$:
$\mu(\lambda_p(n)) < \mu(n)$ for every node~$n$.
Since $\mathbb{N}$ is well-ordered, there is no infinite
decreasing sequence.  The pipeline applies at most~$P$
phases, each decreasing~$\mu$.  After at most $\mu(n_0)$
total lowering steps (where $n_0$ is the initial IR state),
the measure reaches zero and no further lowering is possible.
Therefore the pipeline terminates regardless of the ordering.
\end{proof}

% ============================================================
%  Section 8 — Discussion
% ============================================================
\section{Discussion}\label{sec:discussion}

\paragraph{Phase dependency constraints.}
Not all phase orderings are valid.  The $\LowHint$ advisor
should suggest only orderings that respect phase dependencies.
In the current implementation, the advisor receives the
dependency graph as part of the prompt and is expected to
respect it.  A more robust approach would be to filter
invalid orderings at the kernel level.

\paragraph{Solver cost of node suggestion validation.}
Each node suggestion requires a solver call to confirm
equisatisfiability.  For large IR graphs, this can be
expensive.  Batching suggestions and checking them in a
single solver invocation is a potential optimisation.

\paragraph{Debugging with IRAssist.}
The IR assist advisor is most useful when a lowering phase
fails unexpectedly.  It analyses the IR state and suggests
corrective actions, such as reverting a node rewrite or
applying a different phase.  The trust ceiling ensures that
these suggestions are not automatically applied.

\paragraph{Comparison with compiler IR passes.}
The IR advisory pipeline is analogous to a compiler pass
manager, but with trust-calibrated guidance.  Unlike
compilers, where pass ordering is typically fixed, the
JuGeo pipeline allows dynamic reordering based on the
current IR state.

\paragraph{Integration with lifecycle manager.}
The IR advisors operate during the \texttt{RUNNING} phase
of the kernel lifecycle.  They can be paused during
\texttt{DRAINING} and are inactive during
\texttt{SHUTTING\_DOWN}.

% ============================================================
%  Section 9 — Related Work
% ============================================================
\section{Related Work}\label{sec:related}

\paragraph{Compiler IR optimisation.}
LLVM~\cite{lattner2004llvm} and GCC use pass managers to
order IR transformations.  MLIR~\cite{lattner2020mlir}
provides a multi-level IR infrastructure.  Our work
applies similar ideas to verification IR with
trust calibration.

\paragraph{AI-guided compilation.}
AutoPhase~\cite{huang2019autophase} uses reinforcement
learning for compiler phase ordering.  CompilerGym~\cite{cummins2022compilergym}
provides an RL environment for compiler optimisation.
Our $\LowHint$ advisor uses a language model instead of RL.

\paragraph{Proof term optimisation.}
\LEAN{}~4~\cite{moura2021lean4} performs proof term
elaboration with various reduction strategies.
Coq~\cite{bertot2004coq} supports conversion rules.
Our advisors suggest optimisations rather than
performing them directly.

\paragraph{Copilot advisors.}
Papers~80--83 of this series introduce advisors for
heap, scope, import, callable, research, experiment,
and optimisation tasks.  The IR advisors complement
these with low-level guidance.

\paragraph{Program synthesis for IR.}
Superoptimisation~\cite{massalin1987superoptimizer}
searches for optimal instruction sequences.
Denali~\cite{joshi2002denali} uses theorem proving
for superoptimisation.  Our node suggestor is
analogous but operates on verification IR rather
than machine code.

% ============================================================
%  Section 10 — Conclusion
% ============================================================
\section{Conclusion}\label{sec:conclusion}

We have introduced three Copilot-guided advisors for
IR-level workflows in the \JuGeo{} framework:
\begin{itemize}
\item $\LowHint$ for lowering strategy suggestion and
      phase ordering,
\item $\NodeSugg$ for IR node selection and transformation,
\item $\IRAssist$ for IR-level debugging and optimisation.
\end{itemize}
Five theorems establish correctness preservation, node
suggestion soundness, trust ceiling enforcement,
composition associativity, and phase ordering convergence.
The advisors compose into a pipeline architecture that
integrates with the channel pool and router.

Future work includes learned phase ordering via
reinforcement learning, batched solver validation,
and extension to multi-level IR architectures.

% ============================================================
%  Bibliography
% ============================================================
\begin{thebibliography}{25}

\bibitem{lattner2004llvm}
C.~Lattner and V.~Adve,
``LLVM: a compilation framework for lifelong program analysis
and transformation,''
\emph{CGO}, 2004.

\bibitem{lattner2020mlir}
C.~Lattner \etal{},
``MLIR: scaling compiler infrastructure for domain-specific
computation,''
\emph{CGO}, 2021.

\bibitem{huang2019autophase}
Q.~Huang, A.~Haj-Ali, W.~Moses, J.~Xiang, I.~Stoica,
K.~Asanovic, and J.~Wawrzynek,
``AutoPhase: juggling HLS phase orderings in random forests
with deep reinforcement learning,''
\emph{MLSys}, 2020.

\bibitem{cummins2022compilergym}
C.~Cummins \etal{},
``CompilerGym: robust, performant compiler optimization
environments for AI research,''
\emph{CGO}, 2022.

\bibitem{moura2021lean4}
L.~de~Moura and S.~Ullrich,
``The Lean~4 theorem prover and programming language,''
\emph{CADE}, 2021.

\bibitem{bertot2004coq}
Y.~Bertot and P.~Cast\'{e}ran,
\emph{Interactive Theorem Proving and Program Development:
Coq'Art}.
Springer, 2004.

\bibitem{massalin1987superoptimizer}
H.~Massalin,
``Superoptimizer: a look at the smallest program,''
\emph{ASPLOS}, 1987.

\bibitem{joshi2002denali}
R.~Joshi, G.~Nelson, and K.~Randall,
``Denali: a goal-directed superoptimizer,''
\emph{PLDI}, 2002.

\bibitem{young2024jugeo}
H.~Young,
``Judgment geometry: a framework for trust-calibrated
verification,''
\emph{JuGeo Technical Reports}, Paper~1, 2024.

\bibitem{young2024channels}
H.~Young,
``Evidence channels in judgment geometry,''
\emph{JuGeo Technical Reports}, Paper~88, 2024.

\bibitem{young2024advisors}
H.~Young,
``Copilot advisors for heap, scope, import, and callable
analysis,''
\emph{JuGeo Technical Reports}, Papers~80--81, 2024.

\bibitem{young2024research}
H.~Young,
``Copilot research and experiment advisors,''
\emph{JuGeo Technical Reports}, Paper~82, 2024.

\bibitem{young2024optimization}
H.~Young,
``Copilot optimization advisor,''
\emph{JuGeo Technical Reports}, Paper~83, 2024.

\bibitem{young2024lifecycle}
H.~Young,
``Kernel lifecycle and health monitoring in judgment geometry,''
\emph{JuGeo Technical Reports}, Paper~93, 2024.

\bibitem{demoura2008z3}
L.~de~Moura and N.~Bj{\o}rner,
``Z3: an efficient SMT solver,''
\emph{TACAS}, 2008.

\bibitem{polu2020gptf}
S.~Polu and I.~Sutskever,
``Generative language modeling for automated theorem proving,''
\emph{arXiv preprint arXiv:2009.03393}, 2020.

\bibitem{lample2022hypertree}
G.~Lample \etal{},
``HyperTree proof search for neural theorem proving,''
\emph{NeurIPS}, 2022.

\bibitem{first2022diversity}
E.~First, M.~Rabe, T.~Ringer, and Y.~Brun,
``Diversity-driven automated formal verification,''
\emph{ICSE}, 2022.

\bibitem{yang2019learning}
K.~Yang and J.~Deng,
``Learning to prove theorems via interacting with proof
assistants,''
\emph{ICML}, 2019.

\bibitem{jiang2022thor}
A.~Q.~Jiang, W.~Li, S.~Tworber, and Y.~Wu,
``Thor: wielding hammers to integrate language models and
automated theorem provers,''
\emph{NeurIPS}, 2022.

\end{thebibliography}

% ============================================================
%  Appendix
% ============================================================
\appendix

\section{Auxiliary Definitions and Proofs}\label{app:proofs}

\begin{lemma}[Equisatisfiability is an equivalence relation]%
\label{lem:equisatequiv}
The relation $\equiv_{\mathrm{sat}}$ is reflexive, symmetric,
and transitive.
\end{lemma}
\begin{proof}
\emph{Reflexive:} $n \equiv_{\mathrm{sat}} n$ since the
same formula has the same satisfiability status.

\emph{Symmetric:} if $n \equiv_{\mathrm{sat}} n'$, then
$\varphi(n)$ and $\varphi(n')$ have the same status, so
$n' \equiv_{\mathrm{sat}} n$.

\emph{Transitive:} if $n \equiv_{\mathrm{sat}} n'$ and
$n' \equiv_{\mathrm{sat}} n''$, then all three formulas
have the same status, so $n \equiv_{\mathrm{sat}} n''$.
\end{proof}

\begin{lemma}[Lowering functor preserves identity]%
\label{lem:lowfuncid}
The lowering functor sends the identity morphism $p \to p$
to the identity function $\mathrm{id}_{\mathcal{N}_p}$.
\end{lemma}
\begin{proof}
The identity morphism has no lowering phases to compose.
The empty composition is the identity function.
\end{proof}

\begin{lemma}[Lowering functor preserves composition]%
\label{lem:lowfunccomp}
For morphisms $p \to q$ and $q \to r$ in $\mathbf{Phase}$:
\[
  \Lambda(p \to r) = \Lambda(q \to r) \circ \Lambda(p \to q).
\]
\end{lemma}
\begin{proof}
$\Lambda(p \to r) = \lambda_{r-1} \circ \cdots \circ \lambda_p$.
$\Lambda(q \to r) = \lambda_{r-1} \circ \cdots \circ \lambda_q$.
$\Lambda(p \to q) = \lambda_{q-1} \circ \cdots \circ \lambda_p$.
Their composition is
$\lambda_{r-1} \circ \cdots \circ \lambda_q \circ
\lambda_{q-1} \circ \cdots \circ \lambda_p = \Lambda(p \to r)$.
\end{proof}

\begin{corollary}[$\Lambda$ is a functor]\label{cor:lambdafunc}
$\Lambda \colon \mathbf{Phase} \to \mathbf{Set}$ is a covariant
functor.
\end{corollary}
\begin{proof}
$\Lambda$ preserves identities (Lemma~\ref{lem:lowfuncid})
and composition (Lemma~\ref{lem:lowfunccomp}).
\end{proof}

\begin{lemma}[Valid rewrite rules form a monoid]%
\label{lem:rewritemonoid}
The set of valid rewrite rules on $\mathcal{N}_p$ forms a
monoid under composition, with the identity rewrite as unit.
\end{lemma}
\begin{proof}
\emph{Closure:} if $r_1$ and $r_2$ are valid, then for any
node~$n$: $n \equiv_{\mathrm{sat}} r_1(n) \equiv_{\mathrm{sat}}
r_2(r_1(n))$, so $r_2 \circ r_1$ is valid.

\emph{Identity:} the identity function is valid since
$n \equiv_{\mathrm{sat}} n$.

\emph{Associativity:} function composition is associative.
\end{proof}

\begin{lemma}[Pipeline trust is the minimum]%
\label{lem:pipetrust}
For any pipeline $(A_1, \ldots, A_k)$ of IR advisors, the
pipeline trust is $\min(\tau_1, \ldots, \tau_k)$.
\end{lemma}
\begin{proof}
By induction on~$k$.  For $k = 1$, the pipeline trust is
$\tau_1 = \min(\tau_1)$.  For the inductive step, the
pipeline $(A_1, \ldots, A_k)$ has trust
$\min(\mathrm{trust}(A_1 \mathbin{;} \cdots \mathbin{;}
A_{k-1}), \tau_k) = \min(\min(\tau_1, \ldots, \tau_{k-1}),
\tau_k) = \min(\tau_1, \ldots, \tau_k)$.
\end{proof}

\begin{lemma}[Phase ordering does not affect final result]%
\label{lem:phaseorderindep}
If all lowering phases commute pairwise (\ie{}
$\lambda_p \circ \lambda_q = \lambda_q \circ \lambda_p$
for all $p, q$), then the final lowered IR is independent
of phase ordering.
\end{lemma}
\begin{proof}
Any two orderings of commuting functions yield the same
composed function.  Formally, if $\sigma$ and $\sigma'$
are two permutations, then
$\lambda_{\sigma(P)} \circ \cdots \circ \lambda_{\sigma(1)}
= \lambda_{\sigma'(P)} \circ \cdots \circ \lambda_{\sigma'(1)}$
by iterated commutativity.
\end{proof}

\begin{remark}
In general, lowering phases do not commute.  The
$\LowHint$ advisor's value lies precisely in selecting
a good ordering when phases do not commute.
\end{remark}

\begin{lemma}[Node suggestion confidence calibration]%
\label{lem:confcalib}
If the node suggestor's confidence scores are calibrated
(\ie{} a confidence of~$c$ means the suggestion is correct
with probability~$c$), then accepting only suggestions with
$c \geq 0.9$ yields a 90\% precision rate.
\end{lemma}
\begin{proof}
By the calibration assumption, the probability that a
suggestion with confidence~$c$ is correct equals~$c$.
For $c \geq 0.9$, this probability is at least 0.9.
Therefore, among accepted suggestions, at least 90\%
are correct.
\end{proof}

\begin{lemma}[IR graph size bound under rewrites]%
\label{lem:irgraphbound}
If every rewrite rule is node-local (\ie{} it transforms
one node without adding or removing nodes), then the IR
graph size is invariant under rewrites.
\end{lemma}
\begin{proof}
A node-local rewrite replaces one node with one node.
The number of nodes is unchanged.  The graph structure
(edges) may change, but the vertex count is constant.
\end{proof}

\begin{proposition}[Full pipeline correctness]%
\label{prop:fullpipecorrect}
If the lowering functor is correct and all applied rewrite
rules are valid, then the final IR state is equisatisfiable
with the initial state.
\end{proposition}
\begin{proof}
The pipeline consists of lowering phases (each correct)
and node rewrites (each valid).  By Theorem~\ref{thm:lowcorrect},
the lowering phases preserve equisatisfiability.  By
Theorem~\ref{thm:nodesound}, the rewrites preserve
equisatisfiability.  By transitivity, the final state
is equisatisfiable with the initial state.
\end{proof}

\begin{remark}[Extension to multi-level IR]
The phase category can be generalised to a partially
ordered set of levels, where each level has its own set
of nodes and lowering functions.  The functor construction
and correctness proofs carry over with the appropriate
modifications.  We leave the formal development to future work.
\end{remark}

\begin{lemma}[IR advisor pipeline is idempotent on normal forms]%
\label{lem:irpipeidempotent}
If the IR state is already in its lowest phase (fully lowered),
then running the IR advisory pipeline produces no changes.
\end{lemma}
\begin{proof}
A fully lowered IR state has no remaining phases.  The
$\LowHint$ advisor receives an empty phase list and returns
no hints.  The $\NodeSugg$ advisor receives a fully simplified
graph and returns no rewrites (assuming normal-form detection).
The $\IRAssist$ advisor produces diagnostics but no transformation
suggestions.  Therefore the pipeline is a no-op.
\end{proof}

\begin{lemma}[Pipeline health check precondition]%
\label{lem:pipehealthpre}
The IR advisory pipeline checks pool health before proceeding.
If the pool is unhealthy, the pipeline returns an error without
modifying the IR state.
\end{lemma}
\begin{proof}
Listing~\ref{lst:irpipeline} shows that the first operation is
\texttt{pool.health\_check\_all()}.  If the result is falsy,
the pipeline returns \texttt{\{"error": "Channel pool unhealthy"\}}
immediately.  No IR transformation is attempted.
\end{proof}

\end{document}
